\chapter{Introduction}
\label{cap:introduction}
\addcontentsline{toc}{chapter}{Introduction}
\ref{cap:introduccion}.

\chapterquote{Cuanto peor, mejor para todos. Y cuanto peor para todos, mejor. Mejor para mí, el suyo. Beneficio político.}{-Rajoy}

\section{Introducción}
This report corresponds to the Bachelor’s Thesis in Video Games by Sheila Julvez and José Moreno. In this project, we focused on studying, designing, and implementing a technological solution that allows the use of \textbf{augmented reality (AR)} and spatial detection techniques in live broadcasts, using the open-source platform \textbf{OBS Studio} as the base. The idea arises in a context where the audiovisual world is changing rapidly, combining real-time content creation with immersive technologies that are transforming the way digital content is communicated and consumed.

The starting point of this work lies in how AR has gone from being an experimental curiosity to becoming a common tool in audiovisual production. A clear example is the data from \textbf{\textbackslash{}cite\{GrandViewAR\}}, which estimated the global AR market at around \textbf{\$70.2 billion} in 2023. Nowadays, it is common to see news programs, sports broadcasts, or entertainment shows including 3D graphics, animations, and interactive elements overlaid on real environments, creating a much more visual and engaging experience for viewers.

With this project, we want to test how feasible it is to bring these ideas into a specific context: \textbf{live streaming programming contests}. Although AR is still not widely used in this type of content, we believe it has great potential to make broadcasts clearer, more dynamic, and entertaining, while also highlighting the value of these technical competitions.

\section{Motivation}

 The choice of this topic for our Bachelor’s Thesis comes from two main reasons: on one hand, our academic and professional interest in immersive technologies applied to the audiovisual field; and on the other hand, the growing need to innovate in how technical content is presented and consumed in digital environments, such as live programming contests.

As Video Game Development students, we have worked extensively with graphics systems, interactive environments, and real-time rendering. All of this has sparked our interest in the intersection between advanced graphics technologies and modern media. Augmented reality, in particular, seems very interesting to us because it has applications beyond interactive entertainment, in contexts where clarity, visual appeal, and technical efficiency are also important.

On a more personal level, this project also gives us the opportunity to put into practice what we have learned during the degree in developing interactive graphical solutions, tackling real technical challenges such as integrating AR libraries, detecting physical space, or inserting real-time graphics into a live streaming workflow with OBS Studio. Being able to apply all this knowledge in a functional scenario, with both technical and creative components, is an exciting challenge and we believe it can make an original contribution to a field that is rapidly growing.

\section{Objectives}

 The main goal of this project is to create a theoretical and practical foundation that allows us to understand and apply \textbf{augmented reality (AR)} in live broadcasts using \textbf{OBS Studio}. We focused mainly on how this technology can be used for streaming programming contests.

Beyond exploring what AR can bring visually or communicatively, we also wanted to analyze the technical challenges that arise when integrating it into a real-time broadcast system. These include, for example, detecting and modeling physical space, tracking cameras and moving objects, and rendering 3D elements live without compromising the quality or fluidity of the stream.

Through this thesis, we aim to establish a starting point that can support future projects in this area. Our intention is to contribute to the field of immersive digital communication and provide a reference for students, professionals, or creators who want to experiment with new ways to enhance the viewer experience in live broadcasts, especially in such a technically specific area as programming contests.

\section{Working Plan}

 The development of this Bachelor’s Thesis will follow a sequential structure divided into four main phases, each addressing a specific set of technical and analytical activities. This methodological organization is intended to ensure a logical and coherent progression toward achieving the objectives, allowing effective management of time and resources throughout the research and development process.

The first phase focuses on a thorough literature review. In this initial stage, academic sources, relevant case studies, and applied examples of augmented reality systems used in live broadcasting will be collected and analyzed. Additionally, a comparative evaluation of available technologies and tools will be carried out to select the most suitable ones for the subsequent prototype implementation.

The second phase deals with the conceptual and technical design of the system, along with the progressive development of a functional prototype. This stage includes defining technical and operational requirements, selecting development environments and integration platforms, and programming the first versions of the AR system adapted to streaming environments. It is considered a testing phase to help familiarize ourselves with the various tools studied previously.

The third phase focuses on the implementation and experimental validation of the prototype under controlled conditions. Performance, stability, and visual quality tests will be carried out, along with user experience evaluations and trials in live programming contests.

Finally, the fourth phase is dedicated to the preparation and formalization of the final thesis document. This includes writing the technical report, organizing the results, reviewing and correcting the content, and preparing materials for the public presentation and academic defense of the project before the evaluation committee.

 








